% ================= 3. PRELIMINARY SIZING =================
\section{Preliminary Sizing}
\label{sec:sizing}

This section establishes, before any part is committed to fabrication, that a
cube built to a candidate mass and edge length can actually perform the
jump-up manoeuvre and carry the components the design assumes. It states what
sizing must satisfy (Section~\ref{sec:requirements}), the physics that turns a
target manoeuvre into a required wheel inertia (Section~\ref{sec:jumpup}), the
packaging constraint that closes the loop between cube size and wheel size
(Section~\ref{sec:massbudget}), the wheel design that meets the resulting
target (Section~\ref{sec:rw_sizing}), and the checks that confirm the closure
is real rather than assumed (Section~\ref{sec:sizing_verification}).

Sizing and the control loop are not a sequence of independent steps: they are
solved as one problem. The wheel inertia fixed here sets the angular momentum
the controller has available to arrest a fall and the torque it has available
to do so quickly, so a wheel sized without reference to the control law can be
mechanically sufficient and dynamically useless. In the other direction, the
control law's own behaviour under saturation --- how aggressively it must
desaturate stored wheel momentum, and how much recovery angle it is asked to
cover --- determines how much margin the sizing needs to carry. Each of the
two is therefore an input the other depends on, and the closure procedure of
Section~\ref{sec:massbudget} is the fixed-point iteration that reconciles
them, rather than a one-pass mechanical calculation handed to the controller
afterwards.

\subsection{Requirements Traceability}
\label{sec:requirements}

Each project objective (Section~\ref{sec:objectives}) imposes a required capability.
The mapping below keeps every requirement traceable from motivation to capability,
and is what the sizing and control work below must be shown to satisfy.

\begin{table}[htbp]
\centering
\caption[Objective-to-requirement traceability]{Objective-to-requirement traceability.}
\label{tab:traceability}
\begin{tabular}{lp{0.65\textwidth}}
\toprule
Objective & Required capability \\
\midrule
Edge balance        & The system must actively maintain dynamic stability and regulate its posture about a single line-contact equilibrium position. \\[0.5em]
Corner balance      & The system must maintain multi-axis closed-loop balance and attitude control around a highly unstable point-contact equilibrium. \\[0.5em]
Controlled rotation & The controller must accurately track dynamic reference trajectories to execute smooth slew maneuvers between distinct equilibrium states. \\[0.5em]
Walking (stretch)   & The system must execute a sequential series of dynamic jump-up maneuvers and rapidly re-stabilize itself upon each landing impact. \\
\bottomrule
\end{tabular}
\end{table}

\subsection{Jump-Up Feasibility Analysis}
\label{sec:jumpup}

Given CAD-derived $I_b, I_w, m_b, l_b, m_w, l$, the wheel speed $\omega_w$
required for jump-up is estimated before fabrication, assuming a perfectly
inelastic wheel-body collision. Conservation of angular momentum at impact and
of energy from impact to edge-standing yield
\begin{equation}
\omega_w^2 = (2-\sqrt2)\,\frac{I_w+I_b+m_wl^2}{I_w^2}\,(m_bl_b+m_wl)\,g.
\end{equation}
This value is checked against the motor's achievable speed at the available
bus voltage and used to size wheel inertia (radius, rim mass) within the
motor's operating range without excessively reducing the balancing recovery
angle. Post-fabrication, $\omega_w$ is refined experimentally to correct for
deviations from the ideal collision assumption.

The braking mechanism uses a servo-driven barrier that strikes a bolt head on
the wheel to stop it abruptly. Servo response time bounds the validity of the
``instantaneous stop'' assumption; a slower stop yields a lower effective
post-impact body velocity than predicted, compensated through the same tuning
process.

\subsubsection{Torque-Limited Wheel-Inertia Target}
\label{sec:jumpup_target}

The relation above assumes an instantaneous, lossless momentum transfer. For
preliminary sizing, the achievable motor speed is treated as fixed and the
required per-wheel inertia is derived from a momentum budget that additionally
accounts for finite brake torque. This gives an ideal per-wheel angular
momentum
\begin{equation}
h_{w,\text{ideal}} = \sqrt{\eta}\;\frac{\sqrt{2 g\,\lambda\,\kappa}}{n}\;M L^{3/2},
\label{eq:hw_ideal}
\end{equation}
where $M, L, g$ are cube mass, edge length, and gravity, $\eta$ is an energy
margin, and $\kappa,\lambda,n$ are dimensionless geometric factors set by the
contact case (Table~\ref{tab:contact_factors}). The corner case is more
demanding and governs the design.

\begin{table}[htbp]
\centering
\caption[Contact-case geometric factors]{Contact-case geometric factors.}
\label{tab:contact_factors}
\begin{tabular}{lccc}
\toprule
Case & $\kappa$ (pivot inertia) & $\lambda$ (CoM rise) & $n$ (momentum transfer) \\
\midrule
Edge & $2/3$ & $1/\sqrt2 - 1/2$ & $1$ \\
Corner (governs design) & $11/12$ & $\sqrt3/2 - 1/2$ & $\sqrt2$ \\
\bottomrule
\end{tabular}
\end{table}

A finite brake torque $\tau_b$ cannot deliver the impulse instantaneously and
must exceed the gravitational tip-over torque $\tau_g = MgL/2$. This loss is
captured by a brake-amplification factor
\begin{equation}
\beta = \frac{\tau_b}{\tau_b - \tau_g},
\qquad
h_w = \beta\,h_{w,\text{ideal}},
\qquad
I_{w,\text{target}} = \frac{h_w}{\omega_{\max}},
\label{eq:Iw_target}
\end{equation}
where $\omega_{\max}$ is the maximum practical wheel speed.
Equations~\eqref{eq:hw_ideal} and \eqref{eq:Iw_target} are the sizing relations
used at step~4 of the closure procedure of Section~\ref{sec:massbudget}: given a
candidate cube mass and edge length they return the inertia the wheel of
Section~\ref{sec:rw_sizing} must meet. Their inputs are recorded in
Table~\ref{tab:jumpup_budget} and are not fixed here---$M$ and $L$ are outputs
of the packaging closure, not assumptions preceding it.

Two features of these relations govern how the iteration behaves. First,
$h_w \propto M L^{3/2}$, so growing the cube to accommodate a larger wheel also
raises the requirement; convergence depends on the achievable inertia growing
faster ($\sim D_w^2$) than the requirement does. Second, $\beta$ diverges as
$\tau_b \to \tau_g$: a brake whose torque only marginally exceeds the
gravitational tip-over torque cannot deliver the required impulse at any wheel
speed. The brake torque must therefore be checked against $\tau_g = MgL/2$ on
every iteration, and a comfortable ratio $\tau_b/\tau_g$ maintained rather than
a merely positive margin.

\begin{table}[htbp]
\centering
\caption[Torque-limited wheel-inertia target (corner case)]{Torque-limited wheel-inertia target (corner case). Inputs are fixed by
the closure procedure of Section~\ref{sec:massbudget}; the table is populated on
each iteration. Values are those of the converged iteration
(\texttt{analysis/SIZING.py}, Algorithm~\ref{alg:sizing_closure}): $M$ is a
fixed point of the closure, not an assumption, reproduced exactly
($\delta=\qty{0.0}{\gram}$) by the Stage~6 roll-up at the rounding used in the
bill of materials.
$\tau_b$ is an assumed value pending measurement, justified below and
equivalent to a \qty{40}{\milli\second} stop at the present $h_w$.}
\label{tab:jumpup_budget}
\begin{tabular}{llr}
\toprule
Symbol & Description & Value \\
\midrule
$M$                & Cube mass                    & \qty{1.4654}{\kilo\gram} \\
$L$                & Cube edge length             & \qty{149}{\milli\metre} \\
$\eta$             & Energy margin                & \num{1.05} \\
$\tau_b$           & Brake torque per wheel       & \qty{5}{\newton\metre} (assumed; see below) \\
$\tau_g = MgL/2$   & Gravitational tip-over floor & \qty{1.071}{\newton\metre} \\
$\beta$            & Brake-amplification factor   & \num{1.273} \\
$\omega_{\max}$    & Max wheel speed (design)     & \qty{6000}{\rpm}\textsuperscript{*} \\
$h_w$              & Required angular momentum    & \qty{0.1994}{\kilo\gram\metre\squared\per\second} \\
$I_{w,\text{target}}$ & \textbf{Target wheel inertia} & \qty{3.1735e-4}{\kilogram\metre\squared} \\
\bottomrule
\end{tabular}

\vspace{0.4em}
{\footnotesize\RaggedRight\textsuperscript{*}\,The sizing has additionally been
run at \qty{5000}{\rpm} as an over-dimensioning check. Since
$I_{w,\text{target}} = h_w/\omega_{\max}$, the lower speed raises the inertia
requirement by a factor $6/5$ and so demands more ballast; a wheel that closes
at \qty{5000}{\rpm} therefore closes at \qty{6000}{\rpm} with margin in hand.
\qty{6000}{\rpm} is the value the design claims, \qty{5000}{\rpm} the
conservative case it was verified against.\par}
\end{table}

\paragraph{Selection of $\omega_{\max}$.} The maximum wheel speed is a design
choice rather than a datasheet figure, and it enters the sizing directly
through $I_{w,\text{target}} = h_w/\omega_{\max}$: a higher assumed speed
permits a proportionally smaller wheel. It is therefore the primary relief
variable when the envelope caps the wheel below the requirement (step~5,
Table~\ref{tab:sizing_procedure}), but it must be set from what the drive can
sustain in service, not from the no-load figure $K_v V_{\text{bus}}$. Three
effects separate the two: winding resistance, which costs speed under the torque
required to spin the wheel up; the reduced phase voltage available under
sinusoidal field-oriented modulation; and the fall of pack voltage through the
discharge curve, which must be taken at its lower end if the manoeuvre is to
remain available on a partly-depleted pack rather than only on a fresh one. The
design value is set below the ceiling those effects imply, leaving margin for
motor tolerance, bearing and aerodynamic drag, and the residual uncertainty in
$M$. The resulting back-EMF should remain a comfortable fraction of the usable
bus voltage, and the electrical frequency $\omega_{\max}p/2\pi$ at $p$ pole
pairs must sit inside the driver's commutation capability---both checks to be
confirmed once the motor and pack are selected.

\paragraph{Selection of $\tau_b$.} The brake torque requires more care than the other inputs, because the
mechanism does not generate it in the way the symbol suggests. The brake is a
servo-driven barrier struck by a bolt head on the wheel: the servo only
\emph{holds} the barrier in the wheel's path, and the wheel's own momentum does
the work of the collision. The servo's stall torque is therefore not $\tau_b$
and need not approach it---the TowerPro MG92B of
Table~\ref{tab:bom-supplied} is rated near \qty{0.29}{\newton\metre}, roughly
one seventeenth of the value used in sizing. The servo must resist the barrier
being deflected, not decelerate the wheel.

Because the stop is a collision rather than a friction clutch, $\tau_b$ is not
an independent property of any component and cannot be read from a datasheet.
It is fixed entirely by how quickly the angular momentum is destroyed,
\begin{equation}
\tau_b = \frac{h_w}{\Delta t},
\label{eq:tau_b_from_dt}
\end{equation}
so quoting a brake torque is equivalent to asserting a stopping time. At the
present operating point ($h_w \approx \qty{0.1994}{\kilo\gram\metre\squared\per
\second}$) the design value $\tau_b = \qty{5}{\newton\metre}$ corresponds to
$\Delta t \approx \qty{40}{\milli\second}$. Stated that way the assumption is
open to inspection: \qty{40}{\milli\second} is characteristic of servo travel
rather than of a rigid impact, and a bolt head striking a printed barrier
should arrest the wheel in a few milliseconds, implying a substantially larger
$\tau_b$. The design value is thus expected to be conservative for
\emph{sizing} purposes.

That conservatism is deliberate, and it is asymmetric. Since
$\beta = \tau_b/(\tau_b-\tau_g)$ falls towards unity as $\tau_b$ grows,
underestimating the brake torque inflates the inertia requirement---the safe
direction---while overestimating it leaves the wheel undersized. With the
present $\tau_g$, raising $\tau_b$ from \qtyrange{5}{10}{\newton\metre} relaxes
$I_{w,\text{target}}$ by only about \qty{12}{\percent}, whereas lowering it to
\qty{2}{\newton\metre} inflates the requirement by roughly \qty{69}{\percent}
and to \qty{1.5}{\newton\metre} by roughly \qty{175}{\percent}. The penalty for
error is concentrated almost entirely on the low side, and a mid-range value is
therefore adopted in preference to the higher figure a short impact would
justify.

One consequence must be carried into the structural design rather than left
implicit: $\tau_b$ is also a \emph{load}. The sizing value is chosen for the
dynamics, and the same number should not be used to size the barrier, the bolt,
the servo bracket or the wheel spokes, since a genuine few-millisecond stop
would put an order of magnitude more torque through that load path. Until
measurement settles the question, the structure is to be checked against a
short-$\Delta t$ case and the dynamics against
$\tau_b = \qty{5}{\newton\metre}$. The measurement that closes this is the
spin-down test of Section~\ref{sec:sizing_verification}, which should report both the
impulse-equivalent mean $\bar\tau_b = h_w/\Delta t$, which is the sizing
quantity, and the peak torque, which is the structural one; the two may differ
several-fold, and the logging rate must be fast enough to resolve an event of a
few milliseconds without smoothing the peak away.

\subsection{Sizing Constraints and Closure Procedure}
\label{sec:massbudget}

The cube edge length $L$ and the reaction-wheel inertia $I_w$ are not
independent design choices: they are the two ends of a single constraint chain
that begins with what must physically be carried inside the cube. This section
states that chain and the order in which it is closed. No dimensions are
committed here---component selection is still open (Sections~\ref{sec:cad}
and~\ref{sec:electrical}), and the numbers follow from the procedure rather
than preceding it.

The chain runs in the order set out in Table~\ref{tab:sizing_procedure}. The
volume claim inside the cube is dominated by the LiPo pack and the inner
structure that carries it --- a rigid subframe locating the three motors on
mutually perpendicular axes through a common centre, together with the drivers,
flight controller, IMU and brake servos. The pack is incompressible and sits
near the geometric centre so as not to skew the centre of mass off the body
diagonal, so $L$ is the smallest edge length containing that volume plus wall
thickness and the contact features. $L$ then caps the wheel diameter $D_w$, and
with it the achievable inertia $I_w \sim m_w D_w^2/4$; ballast cannot rescue a
wheel that is short at its maximum diameter, because it necessarily sits inboard
of the ring it displaces.

Against that stands the requirement from the jump-up analysis
(Section~\ref{sec:jumpup_target}),
\[
I_{w,\text{target}} = \frac{h_w(M, L)}{\omega_{\max}},
\qquad
h_w \propto M L^{3/2},
\]
which rises with both $M$ and $L$. Two consequences close the loop. Enlarging
the cube to gain wheel diameter also raises the requirement, but the achievable
inertia grows faster than the requirement does ($D_w^2$ against $L^{3/2}$), so a
wheel short at a given $L$ can be closed by a modest increase in $L$. And any
mass added in doing so feeds back through $h_w \propto M$ and through the
gravitational tip-over torque $\tau_g = MgL/2$ that the brake must exceed ---
which is why the procedure terminates at a fixed point rather than in one pass.

\begin{table}[htbp]
\centering
\caption[Sizing closure procedure]{Sizing closure procedure. Steps are executed in order; a failure at
step 4 returns to step 3 or step 1 as indicated.}
\label{tab:sizing_procedure}
\small
\begin{tabular}{c L{4.6cm} L{6.4cm}}
\toprule
Step & Determines & Governing consideration \\
\midrule
1 & Battery pack geometry & Capacity and discharge rating for the jump-up duty cycle \\
2 & Inner clear volume $V_{\min}$, hence minimum $L$ & Pack $+$ motor subframe $+$ driver/controller stack, non-intersecting; pack near centre \\
3 & Maximum wheel diameter $D_w$, hence achievable $I_w$ & Clearance to subframe, motor body and the two orthogonal modules \\
4 & Requirement $I_{w,\text{target}} = h_w(M,L)/\omega_{\max}$ & Corner jump-up, Eqs.~\eqref{eq:hw_ideal}--\eqref{eq:Iw_target} \\
5 & Accept or iterate & If $I_w < I_{w,\text{target}}$: raise $L$ (step 2), raise $\omega_{\max}$, or trim $M$; re-enter at step 3 \\
6 & Mass reconciliation & Updated $M$ feeds back into step 4; iterate to a fixed point \\
\bottomrule
\end{tabular}
\end{table}

Algorithm~\ref{alg:sizing_closure} (Section~\ref{sec:sizing_algorithm}) states
this procedure formally, as implemented in \texttt{analysis/SIZING.py} and run
to convergence by hand.

The two contact cases are not equally demanding. The corner equilibrium requires a larger rise of the
centre of mass and a less favourable momentum-transfer geometry than the edge
equilibrium, so for a common maximum wheel speed $\omega_{\max}$ it demands the
greater wheel inertia (Table~\ref{tab:contact_factors}). Step~4 above is
therefore always evaluated for the corner case; a wheel that satisfies the
corner requirement satisfies the edge requirement with margin.
Table~\ref{tab:inertia_cases} records both once the procedure closes.

\begin{table}[htbp]
\centering
\caption[Required per-wheel inertia by contact case]{Required per-wheel inertia by contact case, evaluated at
$\omega_{\max}=\qty{6000}{\rpm}$ for the converged $M=\qty{1.4654}{\kilo\gram}$.
Only the geometry factor differs between the cases: $\tau_g$ and $\beta$ are
properties of the cube and the brake, not of the feature it stands on.}
\label{tab:inertia_cases}
\begin{tabular}{lccc}
\toprule
Contact case & $h_w$ (kg\,m$^2$/s) & $I_{w,\text{target}}$ (kg\,m$^2$) & Relative \\
\midrule
Edge                       & \num{0.1809} & \num{2.8790e-4} & \qty{90.7}{\percent} \\
\textbf{Corner (critical)} & \num{0.1994} & \num{3.1735e-4} & \qty{100}{\percent} \\
\bottomrule
\end{tabular}
\end{table}

The corner case demands \qty{9.3}{\percent} more inertia than the edge case, so
sizing on it satisfies the edge requirement with that margin in hand. The
selected wheel delivers \qty{3.6903e-4}{\kilogram\metre\squared}
(Table~\ref{tab:rw_params}), which is \qty{116.3}{\percent} of the corner
requirement and \qty{128.2}{\percent} of the edge requirement.

The allocation of Table~\ref{tab:massbudget} is the bookkeeping side of steps 5
and 6: it is populated as components are selected and re-totalled on each
iteration, since the total is an input to $I_{w,\text{target}}$ as well as an
output of the design.

\begin{table}[htbp]
\centering
\caption[Mass budget (per cube)]{Mass budget (per cube), rolled up from the bill of materials by
\texttt{analysis/SIZING.py} (Stage~6). The total is an input to the inertia
requirement, not only a result: the value below is the \emph{converged} one,
which the closure reproduces exactly ($\delta = \qty{0.0}{\gram}$) against the
$M$ assumed in Table~\ref{tab:jumpup_budget}.}
\label{tab:massbudget}
\small
\begin{tabular}{l c c l}
\toprule
Subsystem & Mass (g) & Share & Basis \\
\midrule
Reaction wheels ($3\times\qty{149.11}{\gram}$) & 447.3 & \qty{30.5}{\percent} & Sized (Sec.~\ref{sec:rw_sizing}) \\
Motors and brake servos ($3\times$ each)      & 245.4 & \qty{16.7}{\percent} & Datasheet \\
Power (battery, XT60, regulator)              & 279.0 & \qty{19.0}{\percent} & Drives envelope (step 1) \\
Electronics (drivers, MCU, encoders, IMU)     &  74.2 & \qty{5.1}{\percent}  & Datasheet \\
Wiring, perfboard and passives                &  45.5 & \qty{3.1}{\percent}  & Estimated \\
\midrule
\textbf{Subtotal --- specified hardware}      & \textbf{1091.4} & \qty{74.5}{\percent} & Measured or datasheet \\
Structure (subframe, frame, fasteners)        & 374.0 & \qty{25.5}{\percent} & \textbf{Allowance, pending CAD} \\
\midrule
\textbf{Total $M$}                            & \textbf{1465.4} & & \textbf{Converged fixed point} \\
\bottomrule
\end{tabular}

\vspace{0.4em}
{\footnotesize\RaggedRight The structure line is an \emph{allowance}, not a
measurement, and at roughly \qty{26}{\percent} of the total it is still the
dominant uncertainty in the budget, even after the \qty{610}{\gram} figure of
the previous iteration was revised down to \qty{374}{\gram}
(subframe \qty{153}{\gram}, frame \qty{184}{\gram}, fasteners/harness
\qty{37}{\gram}, held at the same relative split). It is reported separately
from the specified hardware for that reason: the subtotal is what the design
currently knows, and the total is what it currently projects. Closing the CAD
model replaces the allowance with a real figure and re-enters the loop at
step~6.\par}
\end{table}

\subsection{Reaction-Wheel Sizing}
\label{sec:rw_sizing}

The wheel must supply the per-wheel inertia $I_{w,\text{target}}$ demanded by
the corner jump-up (Section~\ref{sec:jumpup_target}), at an outer diameter no
greater than the envelope permits (step~3 of
Table~\ref{tab:sizing_procedure}), while remaining light enough not to dominate
the cube mass---since its own mass feeds back into $I_{w,\text{target}}$
through $M$. These three demands pull against each other, and the wheel
architecture is chosen so that the design can be trimmed against them after
fabrication rather than having to be right the first time.

\paragraph{Architecture: ballasted ring.} Rather than adjusting inertia by
scaling a solid disc, the wheel is built as a printed \emph{ring plus spokes}
whose inertia is then tuned by \emph{ballast}: standard M6 steel bolt-and-nut
sets fitted through round clearance holes drilled \emph{radially} through the
ring's cross-section, from the inner bore surface to the outer rim surface ---
a true ID-to-OD through-hole, not one drilled through the wheel's face. Each
hole's depth is therefore fixed at the ring's radial width, and its centroid
sits at the ring's mean radius $R_{\text{mean}} = (D_w+D_i)/4$ by construction:
unlike a face-drilled hole there is no free pitch-circle radius to choose. The
hole diameter is instead capped by the wheel's axial thickness, with a printed
wall retained on both faces. Each station both removes plastic and adds steel
at very nearly the largest available radius, so the net inertia gain per
station is large, and the target is reached by choosing how many of the
available stations to populate.

The bolt passing through the ring is also the retention feature: the nut
threads onto it and clamps the ring, so nothing depends on a plastic pocket
floor. Each station carries the centrifugal load
$F = m_{\text{hw}}\,\omega_{\max}^2\,R_{\text{mean}}$ of its own hardware,
reacted by the bolt in tension and by bearing of the head and nut faces ---
now against the ring's curved ID and OD surfaces rather than a flat face, so
both are spot-faced flat (or fitted with a curved washer) to seat squarely ---
checked in Section~\ref{sec:sizing_verification}, and a nyloc or threadlocked
nut is specified so that vibration cannot back it off.

The geometry is defined by the quantities of Table~\ref{tab:rw_params}, all of
which follow from the envelope closure of Section~\ref{sec:massbudget} and are
recorded there once fixed.

\begin{table}[htbp]
\centering
\caption[Reaction-wheel geometry and ballast parameters]{Reaction-wheel
geometry and ballast parameters. Values are the current design point of the
iteration described below; they are provisional while the structural mass of
Table~\ref{tab:massbudget} remains an allowance.}
\label{tab:rw_params}
\small
\begin{tabular}{l L{6.4cm} c}
\toprule
Symbol & Parameter & Value \\
\midrule
$D_w$      & Ring outer diameter (fixed by envelope)   & \qty{120}{\milli\metre} \\
$D_i$      & Ring inner diameter ($D_w - 2w_r$)        & \qty{100}{\milli\metre} \\
$w_r$      & Ring radial width (= ballast hole depth)  & \qty{10}{\milli\metre} \\
$t$        & Wheel axial thickness (caps hole diameter) & \qty{20}{\milli\metre} \\
$n_s, w_s$ & Spoke count and width                     & $3$, \qty{10}{\milli\metre} \\
$\rho_p$   & Printed-material density (PET-CF)          & \qty{1290}{\kilo\gram\per\metre\cubed} \\
$d_h$      & Ballast hole diameter (M6 clearance)      & \qty{6.4}{\milli\metre} \\
$R_{\text{mean}}$ & Ballast centroid radius (forced, $=(D_w{+}D_i)/4$) & \qty{55.0}{\milli\metre} \\
$m_{\text{hw}}$ & Hardware mass per station (bolt $+$ nut) & \qty{7.5}{\gram} \\
$m_b$      & \emph{Net} mass added per station          & \qty{7.085}{\gram} \\
$\Delta I_b$ & \emph{Net} inertia added per station     & \qty{2.151e-5}{\kilogram\metre\squared} \\
$N$        & Station count reaching $I_{w,\text{target}}$ (of $N_{\max}=37$) & \num{3} \\
$m_w$      & Resulting wheel mass                      & \qty{149.11}{\gram} \\
$I_w$      & Resulting wheel inertia                   & \qty{3.6903e-4}{\kilogram\metre\squared} \\
           & \quad as a fraction of $I_{w,\text{target}}$ & \qty{116.3}{\percent} \\
\bottomrule
\end{tabular}
\end{table}

The ballast is characterised \emph{net} of the plastic each hole displaces:
$m_b$ and $\Delta I_b$ are the increments the installed bolt-and-nut adds over
the radial cylinder of material it replaces, evaluated at the forced centroid
radius $R_{\text{mean}}$. Both corrections are taken about the spin axis by
the parallel-axis theorem, retaining each feature's own-axis term rather than
treating it as a point mass --- here the \emph{transverse} solid-cylinder term
$r_h^2/4 + w_r^2/12$, since the hole's own axis is radial and therefore
perpendicular to the spin axis, not parallel to it as a face-drilled hole's
would be. Sizing then reduces to solving
\[
I_{w}(D_w, N) \;=\; I_{w,\text{ring}}(D_w) + N\,\Delta I_b \;\geq\; I_{w,\text{target}}
\]
for the smallest admissible $N$ at each candidate $D_w$, subject to the fit and
mass checks below.

\paragraph{Sizing calculator and iteration procedure.} The relations
above are implemented as a parametric calculator that maps the geometry of
Table~\ref{tab:rw_params} onto the mass and inertia of one wheel, sweeping
candidate outer diameters and returning, for each, the station count required,
the count the gap rule admits, the resulting mass and inertia, and the
pass\slash fail state of every constraint. It is not an optimiser, and the
design point it reports was not found by one.

The outer diameter is not among the quantities it is free to choose. $D_w$ is
\emph{imposed} by the packaging envelope---three orthogonal wheel modules and
the battery must fit within $L$---so the sweep is retained for what it reveals
rather than for what it selects. Left free to minimise wheel mass it prefers a
somewhat larger ring, since $I_{zz}$ scales as $mR^2$ and mass placed at a
large radius buys inertia more cheaply; the free optimum at the present target
is \qty{130}{\milli\metre} (\qty{139.83}{\gram}/wheel, no ballast needed), so at
\qty{120}{\milli\metre} the envelope constraint costs about \qty{9.3}{\gram}
per wheel (\qty{27.8}{\gram} over three) against that unconstrained optimum ---
a modest penalty at the present, lighter design point. This is the price of a
wheel that fits, and a wheel that does not fit is not a candidate.

The procedure actually followed is the fixed-point iteration of
Table~\ref{tab:sizing_procedure}, executed by hand. A cube mass is estimated
and a mass budget assembled from it; the jump-up analysis converts that mass
into an inertia requirement; the wheel is sized against the requirement subject
to the internal-volume and hole-pattern constraints; and the resulting wheel
mass is returned to the budget, which changes the mass that started the loop.
The cycle repeats until it lands on a geometry that simultaneously fits the
internal volume, admits a station count satisfying the inertia requirement, and
reproduces the cube mass it assumed. Convergence is what closes the design:
the reported point is one at which the assumed and the totalled mass agree, so
step~6 of Table~\ref{tab:sizing_procedure} returns approximately zero.

Selection among admissible diameters is made on minimum \emph{wheel mass}
rather than minimum diameter.

\subsection{Sizing Verification}
\label{sec:sizing_verification}

The purpose of this stage is to establish, before any part is committed to
fabrication, that the wheel about to be built can deliver the jump-up manoeuvre.
This is a stricter question than whether the wheel meets its inertia target,
because the target is itself the output of a model whose assumptions are not
exact.

\subsubsection{Two Independent Routes, Compared}

Section~\ref{sec:jumpup} contains two derivations of the jump-up requirement that
proceed from different assumptions, and this stage evaluates both rather than
selecting one. The collision relation of Section~\ref{sec:jumpup} treats the
wheel--body interaction as a perfectly inelastic impact and returns the required
wheel speed $\omega_w$ directly from the body inertia and mass distribution. The
torque-limited momentum budget of Equations~\eqref{eq:hw_ideal}
and~\eqref{eq:Iw_target} instead treats the achievable wheel speed as fixed and
returns the required inertia $I_{w,\text{target}}$, while additionally accounting
for a finite brake torque that the collision relation idealises away.

Agreement between the two is the evidence that the sizing is sound; disagreement
is a finding in its own right and is reported rather than reconciled by choosing
the more favourable result. The comparison is recorded in
Table~\ref{tab:jumpup_gate}.

\begin{table}[htbp]
\centering
\caption[Pre-fabrication feasibility gate]{Pre-fabrication feasibility gate. Pass conditions are stated as
relations because the quantities they relate are outputs of the closure
procedure of Section~\ref{sec:massbudget}, not inputs fixed here.}
\label{tab:jumpup_gate}
\small
\setlength{\tabcolsep}{4pt}
\begin{tabular}{L{4.0cm} L{3.0cm} L{2.4cm} L{3.7cm}}
\toprule
Quantity & Source & Value & Pass condition \\
\midrule
Required wheel speed $\omega_w$ & Collision relation, Sec.~\ref{sec:jumpup} & \textcolor{red}{TBD}\textsuperscript{\textdagger} & $\leq$ achievable speed with margin \\
Achievable wheel speed $\omega_{\max}$ & Drive analysis, Sec.~\ref{sec:jumpup_target} & \qty{6000}{\rpm} & Sustained in service, not no-load \\
Speed margin $\omega_w/\omega_{\max}$ & Derived & \textcolor{red}{TBD}\textsuperscript{\textdagger} & $<1$ by a stated margin \\
Target inertia $I_{w,\text{target}}$ & Eq.~\eqref{eq:Iw_target} & \qty{3.1735e-4}{\kilogram\metre\squared} & --- \\
Achieved inertia $I_w$ & Wheel design, Sec.~\ref{sec:rw_sizing} & \qty{3.6903e-4}{\kilogram\metre\squared} & $\geq I_{w,\text{target}}$ \\
Route disagreement & Both of the above & \qty{+16.3}{\percent} & Reported; explained if large \\
Ballast trim authority & Sec.~\ref{sec:rw_sizing} & $-2 / +34$ stations & Non-zero in both directions \\
Ballast retention margin & Sec.~\ref{sec:rw_sizing} & \qty{1.28}{\percent} of proof; \qty{3.51}{\mega\pascal} bearing & Bolt tension and head/nut bearing stress within allowables \\
\bottomrule
\end{tabular}

\vspace{0.4em}
{\footnotesize\RaggedRight\textsuperscript{\textdagger}\,The collision-route cross-check ($\omega_w$ and the
speed margin) draws on CAD-derived body inertia and mass-distribution figures
($I_b, m_b, l_b$) that are independent of the mass--inertia closure of
Section~\ref{sec:massbudget} and were not re-run as part of this convergence
pass; the values above are cleared pending that re-derivation against the
revised $M$, rather than carried forward from the pre-convergence design
point.\par}
\end{table}

The gate passes on every row except the two pending re-derivation. Trim
authority is now considerably more comfortable in both directions than the
previous iteration. Upward trim is ample: the gap rule admits \num{34} further
stations before $N_{\max}=37$ is reached, enough to more than treble the
inertia. Downward trim is real but tight in absolute terms, because only three
stations are installed to begin with: the wheel sits \qty{16.3}{\percent}
above target while each station is worth roughly \qty{6.8}{\percent} of it, so
two of the three stations may be removed with margin in hand, but removing all
three (returning to the bare ring, which alone reaches only
\qty{95.9}{\percent} of target) fails the requirement. The practical
consequence is unchanged in direction: this design absorbs a
\emph{heavier}-than-predicted cube comfortably and has less room for a lighter
one, which is the favourable direction given that roughly \qty{26}{\percent}
of Table~\ref{tab:massbudget} is still allowance and allowances more often
prove optimistic than pessimistic. Should the structure come in light, the
relief variable is $\omega_{\max}$ rather than the ballast count.

\subsubsection{Brake-Torque Sensitivity}

The brake-amplification factor $\beta = \tau_b/(\tau_b - \tau_g)$ of
Equation~\eqref{eq:Iw_target} diverges as $\tau_b \to \tau_g$, so the inertia
target is not merely uncertain in $\tau_b$ but unboundedly sensitive to it near
the floor. A single assumed value of $\tau_b$ would therefore carry an unstated
and potentially unbounded error into every downstream number.

This stage consequently requires $\tau_b$ to be reported as a sensitivity sweep
rather than as a point value: $I_{w,\text{target}}$ is tabulated across a range
of candidate brake torques, and the achieved wheel inertia is checked against
each. The output of interest is the lowest brake torque at which the design still
closes, which converts an unmeasured parameter into a stated requirement on the
brake. That requirement is discharged when the brake torque is measured on
hardware during the electrical checks of Section~\ref{sec:verification}; until
then the sweep, not an assumption, is what the design rests on.

\subsubsection{Trim Authority}

Because $h_w \propto M L^{3/2}$, a mass overrun raises the inertia requirement
rather than leaving it fixed, and the wheel must therefore retain the ability to
be corrected in both directions after the cube is weighed. The ballast scheme of
Section~\ref{sec:rw_sizing} provides this as a discrete trim, and this stage
records how much of it remains unspent at the selected operating point: a design
sitting only marginally above its target has no authority to absorb a mass
overrun and is recorded as failing this gate even when the nominal inertia check
passes. The trim is reported as the number of ballast elements that may be added
or removed before the target is missed.
